\documentclass[a4paper]{article}
\usepackage{amsfonts}
\usepackage{amsmath}
\usepackage{amssymb}
\usepackage{graphicx}     

\begin{document}
  
\noindent \large Auteur: Abdellah El Moussaoui \\
\noindent \large Studierichting: Informatica\\
\noindent \large Oefeningenreeks 6A nr. 10.6\\

\medskip

\normalsize

Laat de drie termen zijn: $x$, $xq$, $xq^2$\\

$\left\{
\begin{array}{l}
x+xq+xq^2 = 39\\
x^2+\left(xq\right)^2+\left(xq^2\right)^2 = 819\\
\end{array}
\right.$\\

$\Rightarrow x\left(1+q+q^2\right) = 39$\\

$\Rightarrow x=\dfrac{39}{1+q+q^2}$\\

$x^2+x^2q^2+x^2q^4=819$\\

$\Rightarrow x^2\left(1+q^2+q^4\right)=819$\\

$\Rightarrow \left(\dfrac{39}{1+q+q^2}\right)^2 \left(1+q^2+q^4\right)=819$\\

$\Rightarrow \dfrac{1521 \left(1+q^2+q^4\right)}{\left(1+q+q^2\right)^2}=819$\\

$\Rightarrow 1521 \left(1+q^2+q^4\right)=819\left(1+q+q^2\right)^2$\\

$\Rightarrow 13 \left(1+q^2+q^4\right)=7 \left(1+q+q^2\right)^2$\\

$\Rightarrow 13+13q^2+13q^4=7+14q+21q^2+14q^3+7q^4$\\

$\Rightarrow 6q^4-14q^3-8q^2-14q+6=0$\\

$\Rightarrow 3q^4-7q^3-4q^2-7q+3=0$\\

$\Rightarrow \left(q-3\right)\left(3q-1\right) \left(q^2+q+1\right)=0$\\

$\Rightarrow q=3,\dfrac{1}{3}$\\

Voor $q = 3$:\\

$x\left(1+3+9\right)=39 \Rightarrow 13x=39 \Rightarrow x=3$\\

Dus de drie termen zijn:\\

$(3,9,27)$\\

Voor $q = \dfrac{1}{3}$:\\

$x\left(1+\dfrac{1}{3}+\dfrac{1}{9}\right)=39 \Rightarrow x \cdot \dfrac{13}{9}=39 \Rightarrow x=27$\\

Dus de drie termen zijn:\\

$(27,9,3)$\\


\begin{thebibliography}{999}
%\bibitem{a} Referentie
\end{thebibliography}
\end{document} 